% 06 · DBSCAN — density, and the trouble with one threshold
\section{DBSCAN: FOLLOWING THE DENSITY}
\label{sec:dbscan}

\deckslides{slides 13--17, with slide 21 for the single-threshold problem.}

K-means asks how far a star is from a centre. DBSCAN --- density-based
spatial clustering of applications with noise \citep{Ester:96} --- asks a
different question: how many neighbours does a star have within a fixed
radius? That single change buys three things K-means cannot offer: clusters of
arbitrary shape, an explicit notion of noise, and no need to know $K$ in
advance.

\subsection{Definitions}

Fix two parameters: a radius $\varepsilon$ and a minimum count
\texttt{minPts}. Let $N_\varepsilon(\mathbf{x})$ be the set of points within
distance $\varepsilon$ of $\mathbf{x}$ (its $\varepsilon$-neighbourhood).

\begin{itemize}
\item \textbf{Core point.} $\mathbf{x}$ is a core point if
  $|N_\varepsilon(\mathbf{x})| \geq \texttt{minPts}$. Only core points can
  extend a cluster.
\item \textbf{Border point.} A non-core point that lies inside the
  $\varepsilon$-neighbourhood of a core point. It joins that cluster, but
  cannot recruit further points.
\item \textbf{Noise.} Everything else: not core, and not within $\varepsilon$
  of a core point. It is labelled $-1$ and never counted as a member of
  anything.
\item \textbf{Directly density-reachable.} $\mathbf{y}$ is directly
  density-reachable from $\mathbf{x}$ if $\mathbf{y} \in
  N_\varepsilon(\mathbf{x})$ and $\mathbf{x}$ is a core point. Note the
  asymmetry: a border point is reachable \emph{from} its core neighbour but
  not the other way round.
\item \textbf{Density-connected.} Two points are density-connected if there is
  a chain of core points, each directly reachable from the previous, touching
  both. Density-connectivity \emph{is} symmetric, and a DBSCAN cluster is a
  maximal set of density-connected points.
\end{itemize}

\begin{marginnote}[]
\entry{$\varepsilon$ (epsilon)}{The neighbourhood radius. It sets the density
threshold: a point is dense if it has \texttt{minPts} neighbours within
$\varepsilon$ of itself.}
\entry{minPts}{Minimum number of points (including the point itself) required
inside the $\varepsilon$-ball for a point to be a core point. Typical defaults
are $d+1$ for $d$ dimensions, or $\ln n$.}
\end{marginnote}

\begin{figure}[htbp]\centering
  \fig[0.94]{dbscan_film.png}
  \caption{DBSCAN growing a cluster, sampled at four points in the animation:
  core points are identified, a seed point starts a cluster, connected core
  points chain it along the crescent, and border points attach to the chain.
  The points that never enter a core's ball stay noise, permanently.}
  \label{fig:dbscanfim}
\end{figure}

\begin{figure}[htbp]\centering
  \begin{tabular}{@{}c@{\hskip6pt}c@{}}
    \fig[0.44]{dbscan_step1.png} &
    \fig[0.44]{dbscan_step2.png} \\[3pt]
    \fig[0.44]{dbscan_step3.png} &
    \fig[0.44]{dbscan_step4.png}
  \end{tabular}
  \caption{The algorithm in four panels: \panel{a} the two knobs, drawn as
  circles of radius $\varepsilon$ with a count of neighbours --- one point
  passes the \texttt{minPts} test, the other does not; \panel{b} every point
  classified as core (circles), border (triangles) or noise (crosses);
  \panel{c} core points within $\varepsilon$ of each other are chained
  (double-headed arrows), border points hang off the chain (single-headed);
  \panel{d} the result on two interleaved half-moons with uniform noise ---
  both crescents recovered whole, the scattered points discarded.}
  \label{fig:dbscansteps}
\end{figure}

\subsection{The algorithm}

The plain version is short enough to write out. For each point not yet
visited: if it is not a core point, mark it noise for now and move on; if it
is, start a new cluster and grow it by breadth-first search --- every core
point inside the current frontier's $\varepsilon$-neighbourhood joins the
cluster and becomes part of the frontier, and border points are attached as
they are met. The growth stops when the frontier is exhausted. The result does
not depend on the order in which points are visited, except for the assignment
of border points that are within $\varepsilon$ of two different clusters --- an
ambiguity the original paper leaves to the implementation, and a real
reproducibility hazard when comparing libraries.

Cost: with a spatial index, $O(n \log n)$ on average; without one, $O(n^2)$.
The clustering itself is single-pass and cheap; the work is in the
neighbourhood queries, which is the $k$NN machinery of \S\ref{sec:knn} again.

\subsection{Choosing the two knobs}

\texttt{minPts} is the easier of the two. It sets the order of the smallest
structure you are willing to call a cluster, and a common rule of thumb is
\texttt{minPts} $\approx d+1$ for $d$-dimensional data (here 17) or
$\ln n$ for large samples. Our pipeline uses \texttt{minPts} $=5$, which is
deliberately permissive: it keeps small clusters such as the Pleiades
candidates in play, at the cost of admitting noise-driven groups.

$\varepsilon$ is the harder one, and the standard data-driven recipe is a
$k$-distance plot: sort every point's distance to its $k$-th nearest
neighbour, plot the sorted values, and look for the knee --- the point where
the curve leaves the bulk of the data and starts to climb. The knee is where
the density threshold should sit. It is a heuristic, and it is exactly the
step that fails in this project:

\begin{issues}[THE SINGLE-THRESHOLD PROBLEM]
DBSCAN has one density threshold for the whole dataset. Our data have two
populations at very different densities: 357\,056 field stars forming a broad,
dense, chemically smooth background, and 25 clusters of 6--230 members each
sitting inside it. Set $\varepsilon$ for the clusters and the field becomes one
enormous cluster containing everything; set it for the field and the clusters
fragment into noise. There is no middle value that works, and the failure is
not a tuning problem --- it is the assumption of \textbf{Table~\ref{tab:assumptions}}
being false. The next two chapters exist because of it.
\end{issues}

The scale of the effect is visible in the results: on our DR19 run, an
un-normalised abundance space lets HDBSCAN merge the entire field into a
single group (recall $\approx 1.0$ with precision $\approx 0.001$), which is
the signature of a threshold set too high for the background. Row
normalisation (\S\ref{sec:data}) breaks that blob and is the single most
effective precision lever in the pipeline.

\begin{issues}[EXERCISES]
\begin{enumerate}
\item Draw the $\varepsilon$-ball classification by hand for the 2-D dataset
  in \textbf{Figure~\ref{fig:dbscansteps}} with $\varepsilon = 0.35$ and
  \texttt{minPts} $=4$, and mark every point core, border or noise. Then
  change \texttt{minPts} to 6 and say precisely which points change status.
\item Produce the $k$-distance plot for the DR19 member-plus-field matrix at
  $k=5$ and locate the knee. Repeat for $k=2$ and $k=20$. How does the knee
  move, and what does that say about whether a single $\varepsilon$ exists for
  these data?
\item Border points within $\varepsilon$ of two core points of different
  clusters are assigned by visit order in the reference implementation. Find
  two such points in a small synthetic example, and construct two visit orders
  that produce different partitions. Why does this matter when comparing
  scikit-learn with a GPU implementation?
\end{enumerate}
\end{issues}
